\section{Generalización: la limitación fundamental de la IA}

\subsection{La generalización de los modelos está muy por debajo de la humana}

Ilya considera que la generalización es el problema más fundamental de la IA actual:

\begin{importantbox}{Los dos subproblemas de la generalización}
\begin{enumerate}
    \item \textbf{Sample Efficiency}: ¿por qué los modelos necesitan muchos más datos que los seres humanos para aprender lo mismo?
    \item \textbf{Teaching Difficulty}: ¿por qué es mucho más difícil enseñarle algo a un modelo que a un ser humano? Los seres humanos no necesitan verifiable reward: basta con que un mentor muestre su propia forma de pensar para que el estudiante adquiera el método de investigación
\end{enumerate}
\end{importantbox}

\subsection{El papel de la evolución}

Para habilidades como la visión, el oído o el movimiento, la evolución probablemente aportó un prior poderoso (evolutionary prior). La destreza de los dedos humanos supera con mucho a la de los robots, y la capacidad de un niño de 5 años para reconocer automóviles ya basta para la conducción autónoma.

Pero para habilidades \textbf{de aparición reciente}, como el lenguaje, las matemáticas o la programación, es poco probable que la evolución haya aportado priors específicos. Sin embargo, la eficiencia de aprendizaje de los seres humanos en estos campos sigue siendo muy superior a la de los modelos, lo que sugiere que los seres humanos podrían poseer algún \textbf{algoritmo de aprendizaje más fundamental y general}, en lugar de depender solo de priors evolutivos específicos de cada dominio.

\begin{importantbox}{La robustez del aprendizaje humano}
El aprendizaje humano tiene una robustez (robustness) asombrosa:
\begin{itemize}
    \item Menos muestras
    \item Menos supervisión (un teenager no necesita verifiable reward para aprender a manejar)
    \item Una robustez muy fuerte (sigue aprendiendo con eficacia en entornos completamente nuevos)
\end{itemize}
Ilya considera que existe algún \textbf{principio de aprendizaje automático} aún no descubierto que podría lograr una capacidad de generalización semejante: "I think it can be done. The fact that people are like that I think it's a proof that it can be done."
\end{importantbox}

\begin{knowledgebox}{Es posible que las neuronas humanas hagan más cómputo del que se cree}
Ilya menciona un posible blocker: es posible que las neuronas humanas realicen en la práctica más cómputo del que suponemos. Si esto es cierto, y resulta crucial para el aprendizaje, entonces alcanzar una generalización a nivel humano podría requerir más cómputo.
\end{knowledgebox}

\subsection{Resumen de la sección}

La generalización es la limitación más fundamental de la IA actual. Los seres humanos siguen mostrando una eficiencia de aprendizaje excepcional en dominios donde la evolución difícilmente aportó priors, como la programación, lo que sugiere la existencia de algún principio de aprendizaje más fundamental. Ilya tiene ideas al respecto, pero no puede discutirlas públicamente por razones de competencia.


